\section{结果}
\label{zh:sec:results}

\subsection{基准对比与统计评估}
\label{zh:subsec:results_main}

在定理验证方面，$O(k\log k)$ 扫描与穷举枚举在全部 4,000 个实例上完全一致，中心带宽度的最大差值低于 $10^{-12}$。

在主要配对比较中，表~\ref{zh:tab:main_results} 汇总了 3,500 个配对场景的结果。MIKU 将无碰撞到达目标率从基线 B0 的 62.69\% 提高至 77.46\%，实现配对差值 14.77 个百分点，其 95\% Bootstrap 置信区间为 13.54 至 15.97，精确 McNemar 检验 $p=1.94\times10^{-133}$。碰撞率从 2.17\% 降至 1.14\%，降幅达 1.03 个百分点，95\% 置信区间为 $-1.37$ 至 $-0.71$，$p=2.91\times10^{-11}$。平均进度提高 8.93 个百分点，加加速度均方根降低 0.120 m/s$^3$，带惩罚行驶时间减少 0.81 s。平均间隙的总体差值为 $-0.010$ m，95\% 置信区间为 $[-0.024,0.004]$ m。

\begin{table}
\tbl{3,500 个配对场景基准评估结果。比例以百分数表示，连续指标报告均值，运行时间以毫秒表示。}
{%
\begin{tabular}{lrrrr}
\toprule
指标 & B0 & B1 & B2 & MIKU \\
\midrule
无碰撞到达目标（\%） & 62.69 & 52.37 & 62.09 & \textbf{77.46} \\
碰撞（\%） & 2.17 & 2.14 & 2.23 & \textbf{1.14} \\
退化停车（\%） & 35.40 & 45.77 & 35.97 & \textbf{21.51} \\
进度比（\%） & 77.52 & 71.07 & 76.62 & \textbf{86.45} \\
加加速度均方根（m/s$^3$） & 1.252 & 1.315 & 1.300 & \textbf{1.132} \\
运行时间第 50 百分位（ms） & 10.11 & 11.37 & 31.20 & \textbf{9.98} \\
运行时间第 95 百分位（ms） & 56.79 & 63.24 & 169.49 & \textbf{55.93} \\
运行时间第 99 百分位（ms） & 138.88 & 169.73 & 463.05 & 221.80 \\
\bottomrule
\end{tabular}}
\begin{tabnote}[注]
B0：采用贪心逐障碍物边界的 Apollo 基线；B1：未引入到达时间条件的时间展开占用；B2：3 次迭代启发式耦合；MIKU：本文提出的约束编译器。加粗表示各指标最优结果。所有方法均在 7 类基准场景族相同的初始状态与随机种子下配对评测，样本量 $N=3,500$。
\end{tabnote}
\label{zh:tab:main_results}
\end{table}

在类别层面的效应分析中，图~\ref{zh:fig:evidence_dashboard} 面板 (a) 定位了差异的来源：成功率增益最大的是窄间隙多障碍物类别，提升达 65.4 个百分点；其次是延迟交叉类别，提升达 26.8 个百分点。交叉行人和车辆切入类别未出现总体差异，因为其几何在障碍物任一侧都存在一条主导通道；MIKU 的群组级划分与时间同伦在这些单走廊构型下不改变约束接口，证实该方法在基线接口已经充分的场景中不引入性能退化。预测噪声类别的成功率提高 3.2 个百分点，碰撞率降低 7.0 个百分点；然而进度下降 1.12 个百分点，加加速度均方根上升 0.041 m/s$^3$，这是误差界较大时保守的不确定性扩展占用管的预期副效应。面板 (d) 显示配对增益来自退化停车向成功到达的迁移。

在与联合时空格栅搜索参照 B3 的 70 个案例配对对比中，B3 达到 80.0\%，MIKU 为 77.1\%，配对差值为 $-2.86$ 个百分点，95\% 置信区间为 $[-11.43,5.71]$，$p=0.754$；B3 的运行时间中位数为 1,835.43 ms，而 MIKU 为 10.44 ms。

\begin{figure}
\centering
\includegraphics[width=\textwidth]{evidence_dashboard.pdf}
\caption{3,500 个配对场景的结果汇总。各面板依次给出分场景族成功率差值（a）、规划时延分布（b）、消融降幅（c）与结果构成（d）。}
\label{zh:fig:evidence_dashboard}
\end{figure}

\subsection{组件消融与机制归因}
\label{zh:subsec:results_ablation}

在相同的 3,500 个场景上执行减法消融，沿用第~\ref{zh:sec:experimental_design} 节的 5,000 次重采样 Bootstrap。移除 Top-$K$ 空间同伦后成功率下降 10.17 个百分点；移除威胁感知裕度后下降 10.74 个百分点。移除时间同伦图后，延迟交叉类别成功率从 89.6\% 降至 66.8\%，在该类别上造成 22.8 个百分点的损失。移除鲁棒占用管后，预测噪声类别的碰撞率从 3.6\% 上升至 11.0\%。到达时间投影与纵向分组的总体变化较小，其价值集中于特定构型。图~\ref{zh:fig:evidence_dashboard} 面板 (c) 汇总了下表所列的效应。

\begin{table}
\tbl{相同 3,500 个配对场景上的减法消融结果；$\Delta$ 列给出相对完整 MIKU 的带符号变化，以百分点为单位。}
{%
\begin{tabular}{l lrrrrrr}
\toprule
变体 & 移除组件 & 成功率（\%） & $\Delta$ 成功 & 碰撞率（\%） & $\Delta$ 碰撞 & 进度（\%） & $\Delta$ 进度 \\
\midrule
A1 & 到达时间投影 & 77.34 & $-0.12$ & 1.23 & $+0.09$ & 86.73 & $+0.28$ \\
A2 & 纵向分组 & 76.94 & $-0.52$ & 1.14 & $0.00$ & 85.89 & $-0.56$ \\
A3 & Top-$K$ 空间同伦 & 67.29 & $-10.17$ & 1.14 & $0.00$ & 79.53 & $-6.92$ \\
A4 & 威胁感知裕度 & 66.71 & $-10.74$ & 1.14 & $0.00$ & 78.85 & $-7.60$ \\
A5 & 时间同伦图 & 74.20 & $-3.26$ & 1.17 & $+0.03$ & 84.29 & $-2.16$ \\
A6 & 鲁棒占用管 & 76.97 & $-0.49$ & 2.20 & $+1.06$ & 86.60 & $+0.15$ \\
A7 & 交替细化 & 76.63 & $-0.83$ & 1.03 & $-0.11$ & 85.92 & $-0.53$ \\
MIKU & 无 & \textbf{77.46} & --- & \textbf{1.14} & --- & \textbf{86.45} & --- \\
\bottomrule
\end{tabular}}
\begin{tabnote}[注]
基于相同的 $N=3,500$ 场景并在 5,000 次 Bootstrap 重采样下计算。$\Delta$ 表示相对完整 MIKU 流水线的带符号百分点变化；在成功率或进度上的负 $\Delta$ 以及碰撞率上的正 $\Delta$ 代表移除该组件后系统性能发生退化。加粗表示完整 MIKU 框架。
\end{tabnote}
\label{zh:tab:ablation_results}
\end{table}

几何上，如图~\ref{zh:fig:narrow_mechanism_en}所示，在窄通道中，逐障碍物策略不断抬升与压低相互对置的路径边界，直至中心区间被消解为空；而群组级划分在协调障碍物侧向选择后保留一条连续带。在延迟交叉中，时间图同时保留障碍物前通行与障碍物后通行两条候选，并强制其因果顺序。

\begin{figure}
\centering
\includegraphics[width=\textwidth]{fig_narrow_en.pdf}
\caption{贪心逐障碍物边界使紧缩处可行区间为空（a）。MIKU 的群组级划分保留连续中心走廊（b）。}
\label{zh:fig:narrow_mechanism_en}
\end{figure}

\subsection{鲁棒性、时延分布与滚动重规划}
\label{zh:subsec:results_robustness}

在预测噪声子组中，除了第~\ref{zh:subsec:results_main} 节报告的 3.2 个百分点成功率增益与 7.0 个百分点碰撞率降低外，MIKU 的平均进度为 59.73\%，B0 为 60.85\%；加加速度均方根为 2.274 m/s$^3$，B0 为 2.234 m/s$^3$。这与如下解释一致：当误差界较大时，不确定性扩展占用会带来较为保守的行为。

完整调用运行时间的中位数与第 95 百分位与 B0 相当或更低，但第 99 百分位并非如此：MIKU 为 221.80 ms，B0 为 138.88 ms，图~\ref{zh:fig:evidence_dashboard} 面板 (b) 中可见上尾的分化。B2 因始终执行三次协调迭代，在所有分位数上都更慢。

滚动重规划协议评估多周期连续规划下的闭环累积控制表现。在 700 个案例中，成功率从 62.29\% 提升至 72.29\%，实现 10.0 个百分点的改善，95\% 置信区间为 7.57 至 12.57；两种方法的碰撞率均为 1.29\%，平均进度从 82.55\% 升至 90.25\%，单回合运行时间第 95 百分位从 827.20 ms 升至 1,361.51 ms。

\subsection{量产软件接口集成与验证}
\label{zh:subsec:results_apollo}

Apollo 11.0 集成直接通过既有的路径边界与 ST 边界模块传递生成的约束条件，完整保留下游 QP 的目标函数、运动学等式与凸优化结构。如图~\ref{zh:fig:apollo_cases}所示，三个代表性 Dreamview 测试工况，即动态超车、窄通道通行与单车道绕行，在 Apollo 0.5\,s 规划周期内均实现了平滑轨迹生成与稳健避障，证实了方法与量产级自动驾驶系统架构的工程兼容性与可行性。

\begin{figure}
\centering
\begin{minipage}{0.32\textwidth}\centering\includegraphics[width=\linewidth]{dreamview/scn01_during.png}\par\vspace{2pt}\footnotesize (a) 动态超车\end{minipage}\hfill
\begin{minipage}{0.32\textwidth}\centering\includegraphics[width=\linewidth]{dreamview/scn02_during.png}\par\vspace{2pt}\footnotesize (b) 窄通道通行\end{minipage}\hfill
\begin{minipage}{0.32\textwidth}\centering\includegraphics[width=\linewidth]{dreamview/scn03_during.png}\par\vspace{2pt}\footnotesize (c) 单车道绕行\end{minipage}
\caption{Apollo 11.0 Dreamview 中的动态超车（a）、窄通道通行（b）与单车道绕行（c）。}
\label{zh:fig:apollo_cases}
\end{figure}
